\section{Method}
\label{sec:method}

Our method operates in four stages. Stage~1 applies standard geometric post-processing to establish a baseline estimate. Stages~2--4 constitute our contribution: self-conditioned iterative refinement with decoupled rotation-translation fusion.

\subsection{Preliminaries: Geometric Post-Processing (Stage 1)}
\label{sec:stage1}

Given $N_t = 91$ timesteps of $N$-camera images, we run Pi3X to obtain per-timestep poses $\{T_k^c\}_{k=1}^{N_t}$ for each camera $c \in \{1,\ldots,N\}$, where each $T_k^c \in \text{SE}(3)$ is defined up to a per-timestep Sim(3) ambiguity (i.e., each timestep's output coordinate frame is arbitrary in position, orientation, and scale). We aggregate these into a single estimate via three standard steps.

\paragraph{Gauge fixing.} Since each timestep uses an arbitrarily chosen global frame, we align all timesteps to a reference via Sim(3) Procrustes~\cite{umeyama1991}:
\begin{equation}
    \{s^*, R^*, \mathbf{t}^*\} = \argmin_{s, R, \mathbf{t}} \sum_{c=1}^N \| s R \, \mathbf{p}_k^c + \mathbf{t} - \mathbf{p}_{\text{ref}}^c \|^2
\end{equation}
where $\mathbf{p}_k^c \in \mathbb{R}^3$ denotes the position of camera $c$ at timestep $k$, and $s \in \mathbb{R}^+$, $R \in \text{SO}(3)$, $\mathbf{t} \in \mathbb{R}^3$ parametrise the Sim(3) transform.

\paragraph{Outlier rejection.} For each camera, we project its $N_t$ estimated positions onto the dominant 2D plane (via PCA on camera centres) and reject outliers whose in-plane distance exceeds $3\times$ the Median Absolute Deviation (MAD). Inlier rotations are aggregated via the Fr\'{e}chet mean on SO(3)~\cite{moakher2002}---the rotation minimising the sum of squared geodesic distances to all inlier samples.

\paragraph{Pose graph optimization.} We extract pairwise relative poses $T_{ij} = (T^i)^{-1} T^j$ from inlier timesteps, computing the Fr\'{e}chet mean rotation and mean translation for each camera pair $(i, j)$. These relative poses serve as edges in a pose graph, weighted by the inverse variance of the per-timestep estimates (edges with higher consistency receive more weight). We then optimise globally via $L_1$ Weiszfeld rotation averaging on SO(3)~\cite{hartley2013} followed by translation averaging with hard co-location constraints for cameras sharing a physical mounting pole.

\paragraph{Residual error analysis.} This pipeline reduces ATE by 10\% (0.177 $\rightarrow$ 0.161) but leaves ARE unchanged (4.37\degree{} $\rightarrow$ 4.39\degree{}). The persistence of rotation error under robust averaging over 91 independent timesteps indicates that the residual is systematic model bias rather than sample noise. Since the scene provides weak geometric signal (approximately planar static structure, dynamic repetitive content), further classical aggregation cannot improve upon this estimate. This motivates the use of Pi3X's conditioning interface to inject prior geometric information, as developed in the following sections.

\subsection{Self-Conditioned Inference (Stage 2)}
\label{sec:stage2}

Pi3X supports multimodal conditioning: given approximate poses $\hat{T}^c$ and camera intrinsic matrices $\boldsymbol{\Pi}^c \in \mathbb{R}^{3 \times 3}$, it refines its predictions by attending to the provided geometric priors during encoding. We condition on the Stage~1 output:
\begin{equation}
    \{T_k^{c,(1)}\} = f_\theta(\text{imgs}_k \mid \hat{T}^c, \boldsymbol{\Pi}^c), \quad k = 1, \ldots, N_t
\end{equation}
where $f_\theta$ denotes Pi3X and $\hat{T}^c$ is the Stage~1 estimate converted to the model's input camera ordering.

Applying the same geometric post-processing to $\{T_k^{c,(1)}\}$ yields a refined estimate with ARE reduced from $4.39\degree$ to $3.49\degree$ (20\% improvement) but ATE increased from 0.161 to 0.166 (3\% degradation). Iterating this process na\"ively amplifies both effects: rotation improves further (to ${\sim}3.3\degree{}$) while translation degrades monotonically (to 0.195 ATE after 5 iterations). Section~\ref{sec:theory} provides a formal explanation: the rotation component contracts toward a fixed point on SO(3), while translations diverge due to Sim(3) scale mismatch in the conditioning pathway.

\subsection{Decoupled Hybrid Fusion (Stage 3)}
\label{sec:stage3}

Since conditioned inference improves rotation but degrades translation, we decouple them in the pose graph by sourcing edges from different inference runs. Let $R_{ij}$ and $\mathbf{t}_{ij}$ denote the relative rotation and translation from camera $i$ to camera $j$, computed as the mean over inlier timesteps. We construct fused edges:

\begin{equation}
    \tilde{T}_{ij}^{\text{fused}} = \begin{pmatrix} R_{ij}^{\text{cond}} & \mathbf{t}_{ij}^{\text{uncond}} \\ \mathbf{0}^\top & 1 \end{pmatrix}
\end{equation}

where $R_{ij}^{\text{cond}}$ is the mean relative rotation extracted from conditioned inference (Stage~2) and $\mathbf{t}_{ij}^{\text{uncond}}$ is the mean relative translation from the original unconditioned run (Stage~1). Edge weights $w_{ij}$ (inverse variance of the per-timestep relative pose estimates) are taken as $w_{ij} = \max(w_{ij}^{\text{cond}}, w_{ij}^{\text{uncond}})$.

The initial pose for graph optimization similarly uses conditioned rotations with unconditioned translations:
\begin{equation}
    T_{\text{init}}^c = \begin{pmatrix} R_{\text{cond}}^c & \mathbf{t}_{\text{uncond}}^c \\ \mathbf{0}^\top & 1 \end{pmatrix}
\end{equation}

This yields ARE = 3.45\degree{} (rotation improved) with ATE = 0.159 (translation preserved)---capturing the benefits of both sources.

\subsection{Iterated Hybrid Refinement}
\label{sec:stage4}

Since a single conditioned pass reduces rotation error by 20\%, a natural question is whether iterating the process yields further gains. We apply hybrid fusion iteratively: at each step, the optimised rotation estimate from the previous iteration serves as the new conditioning input.

\begin{algorithm}[t]
\caption{Iterated Hybrid Refinement}
\label{alg:itermc}
\begin{algorithmic}[1]
\REQUIRE Pipeline estimate $\hat{T}^c$, images, intrinsics $\boldsymbol{\Pi}^c$
\REQUIRE Iterations: $M$
\STATE $T_{\text{curr}} \leftarrow \hat{T}^c$
\STATE Compute unconditioned edges $\{\mathbf{t}_{ij}^{\text{uncond}}, w_{ij}^t\}$ \COMMENT{fixed throughout}
\FOR{$m = 1$ to $M$}
    \STATE $\{T_{k}^{c,(m)}\} \leftarrow f_\theta(\text{imgs} \mid T_{\text{curr}}, \boldsymbol{\Pi}^c)$ for all $k$
    \STATE Extract rotation edges $\{R_{ij}^{\text{cond}}, w_{ij}^R\}$ from pipeline-processed conditioned poses
    \STATE Fuse: $\tilde{T}_{ij} = (R_{ij}^{\text{cond}}, \mathbf{t}_{ij}^{\text{uncond}})$
    \STATE $T_{\text{opt}} \leftarrow \text{PoseGraph}(\bar{T}_{\text{init}}, \tilde{T}_{ij}, w_{ij})$
    \STATE $T_{\text{curr}} \leftarrow (R_{\text{opt}}, \mathbf{t}_{\text{uncond}})$ \COMMENT{anchor translations}
\ENDFOR
\RETURN $T_{\text{opt}}$
\end{algorithmic}
\end{algorithm}

\paragraph{Design choices.}
\begin{itemize}
    \item \textbf{Translation anchoring}: Translations are \emph{never} updated from conditioned outputs---only from the original unconditioned baseline. This is the key invariant that prevents Sim(3) drift.
    \item \textbf{Early stopping}: Terminate when ARE improvement between iterations falls below $0.1\degree$, or after $M=2$ iterations (empirically sufficient; see Section~\ref{sec:ablation}).
\end{itemize}

\paragraph{Monte Carlo extension.} The iterative procedure admits a natural stochastic extension: perturbing the conditioning poses before each forward pass and pooling the resulting observations across $S$ samples. Under a posterior sampling interpretation (Section~\ref{sec:theory}), this estimates the posterior mean via Monte Carlo integration, with variance decreasing as $O(1/\sqrt{S})$. We implement this with structured perturbation (shared global rotation preserving inter-camera geometry), Sim(3) alignment across samples, and pool-before-pipeline aggregation. However, our experiments (Section~\ref{sec:ablation_S}) demonstrate that this extension provides no benefit: the model's posterior under self-conditioning is so tightly concentrated that $S=1$ (deterministic iteration) is optimal.
